\documentclass[11pt]{article}
\usepackage[margin=1in]{geometry}
\usepackage[T1]{fontenc}
\usepackage{lmodern,microtype}
\usepackage{amsmath,amssymb,amsthm,booktabs}
\usepackage{hyperref}
\hypersetup{hidelinks,pdftitle={},pdfauthor={},pdfsubject={},pdfkeywords={},pdfcreator={},pdfproducer={}}
\pdfinfoomitdate=1
\pdftrailerid{}
\newtheorem{theorem}{Theorem}
\newtheorem{proposition}{Proposition}
\newcommand{\tr}{\operatorname{tr}}
\newcommand{\diag}{\operatorname{diag}}
\newcommand{\dd}{\mathrm d}
\newcommand{\HH}{\mathsf H}
\numberwithin{equation}{section}
\title{A local clock-foliated repair of the inhomogeneous spectrum\\of a rotating Einstein-static scalar reference}
\author{}
\date{29 October 2026}
\begin{document}
\maketitle
\begin{abstract}
An exact rotating eight-scalar configuration on the Einstein cylinder, supplemented by a ninth scalar clock, supplies a proper global spacetime reference map.
A radial potential can make its fixed-charge homogeneous modes elliptic, but this does not control its inhomogeneous degrees of freedom.
We reduce the coupled scalar and vector harmonics, including the metric constraints and exceptional low harmonics.
The canonical model has an oscillatory instability in the lowest scalar harmonic and a growing mode in the third harmonic; the latter persists throughout its homogeneous elliptic parameter window.
We construct a local preferred-clock effective action containing traceless spatial strain, curvature-calibrated compression, and a nonnegative target-alignment potential.
All additions vanish with their first variations on the complete homogeneous ansatz.
At explicit rational parameters, exact principal-minor certificates prove positive kinetic and stiffness forms for every physical inhomogeneous scalar and vector harmonic, with a positive tensor correction.
A stationary pure Gaussian is constructed for a finite scalar band of 141 canonical oscillators, together with a bounded coupled linear metric--matter evolution.
The modified principal cones, cutoff assumption and unchanged homogeneous nonlinear limitations are explicit; neither a continuum Hadamard theorem nor nonlinear stability is asserted.
\end{abstract}

\section{Problem and construction}

An exact reference-field solution, linear spectral stability, positive quadratic energy and nonlinear stability impose different requirements.
This article addresses the inhomogeneous degrees of freedom of a rotating scalar reference on a compact three-sphere.
The exact reference and its radial homogeneous stabilization are starting constructions, rederived here to fix conventions.
The new results are the constrained harmonic reduction, an obstruction to radial repair, and a local effective action that removes the inhomogeneous instabilities while preserving the homogeneous equations.

The target-radial and target-tangent gradient components mix within a scalar representation of the diagonal spatial--target rotation group.
Metric lapse and shift must be eliminated before counting physical modes.
The first scalar harmonic has a residual spatial diffeomorphism, whereas an internal-symmetry zero in the next harmonic is physical.
The sphere-harmonic conventions agree with Ref.~\cite{harmonics}; all embedding identities needed for the reduction are derived below.

The repair uses a timelike clock to define local spatial operators.
Preferred-foliation higher-spatial-derivative theories provide broader context~\cite{horava}, but no ultraviolet-completion or renormalizability result from that literature is assumed.
Our claim concerns a specified local effective action and its constrained quadratic theory.
The central positive statement is Theorem~\ref{thm:positive}.

\section{Canonical reference and homogeneous sector}
\label{sec:background}

Use signature $(-+++)$, $\hbar=c_{\rm light}=1$, $\kappa=8\pi G>0$, eight real canonical scalars $\Phi=(X^A,Z^A)$, $A=1,\ldots,4$, and a canonical clock $\chi$.
Set
\begin{equation}
 s=|X|^2+|Z|^2,\quad b=\frac{2-\alpha}{4},\quad
 V_0(s)=\frac1{\kappa a_0^2}\left[\frac{m}{8b}(\kappa s-b)^2+g_4(\kappa s-b)^4\right].
 \label{eq:potential}
\end{equation}
The action is
\begin{equation}
 S_0=\int\sqrt{-g}\left[\frac{R-2\Lambda}{2\kappa}
 -\frac12\sum_{I=1}^{8}\nabla_\mu\Phi^I\nabla^\mu\Phi^I
 -\frac12\nabla_\mu\chi\nabla^\mu\chi-V_0(s)\right]\dd^4x.
 \label{eq:canonical}
\end{equation}
For $g_4>0$, $V_0$ is bounded below even when $m<0$.
Let $Y:S^3\to\mathbb R^4$ be the unit-sphere embedding, satisfying
\begin{equation}
 Y\cdot Y=1,\quad Y\cdot\dd Y=0,\quad \dd Y\cdot\dd Y=\gamma,
 \quad\nabla_i\nabla_jY=-\gamma_{ij}Y.
\end{equation}
The reference is
\begin{align}
 g_0&=-\dd t^2+a_0^2\gamma,& X+iZ&=r_0e^{i\omega_0t}Y,&\chi&=ct,\label{eq:bg}\\
 r_0^2&=b/\kappa,&\omega_0^2&=3/a_0^2,&
 c^2&=\frac\alpha{\kappa a_0^2},\qquad
 \Lambda=\frac{6+\alpha}{4a_0^2},\quad0<\alpha<2.
\end{align}
At $s_0=r_0^2$, $V_0(s_0)=V_0'(s_0)=0$ and $4a_0^2r_0^2V_0''(s_0)=m$.
Since $\Delta_\gamma Y=-3Y$, every scalar solves its background field equation.
The stress density and pressure are
\begin{equation}
 \rho_0=\frac{6-\alpha}{4\kappa a_0^2},\qquad
 p_0=\frac{2+\alpha}{4\kappa a_0^2},
\end{equation}
and solve the two independent Einstein equations.
The cylinder is complete with compact Cauchy surface.
The nine-field map is a proper global embedding: recover $t=\chi/c$ and
\begin{equation}
 Y=\frac{\cos(\omega_0t)X+\sin(\omega_0t)Z}{r_0}.
\end{equation}
Bounded $\chi$ bounds $t$ and the sphere is compact, proving properness.
The dimensions are $[a_0]=L$, $[\kappa]=L^2$, $[\Phi]=L^{-1}$ and $[V_0]=L^{-4}$.

For the homogeneous ansatz $X+iZ=r(t)e^{i\theta(t)}Y$, the conserved charges are $Q_\theta=a^3r^2\dot\theta$ and $S_\chi=a^3\dot\chi$.
At fixed charges, with $u_a=\delta a/a_0$, $v_r=\delta r/r_0$ and $\tau=t/a_0$, the linear equations are
\begin{equation}
 \begin{pmatrix}u_a\\v_r\end{pmatrix}''=
 \begin{pmatrix}3+\alpha/2&1-\alpha/2\\-12&-12-m\end{pmatrix}
 \begin{pmatrix}u_a\\v_r\end{pmatrix}.
\end{equation}
Two distinct elliptic modes occur for
\begin{gather}
 0<\alpha<4\sqrt{21}-18,\label{eq:windowa}\\
 -15-\alpha/2+\sqrt{24(2-\alpha)}<m<-
 \frac{24(2+\alpha)}{6+\alpha}.\label{eq:windowm}
\end{gather}
The point used here is
\begin{equation}
 \alpha=\tfrac1{10},\quad m=-\tfrac{207}{25},\quad g_4=10,
 \qquad b=\tfrac{19}{40}.
 \label{eq:radialpoint}
\end{equation}
Its homogeneous frequencies are $0.3061068927/a_0$ and $0.7591433134/a_0$.
Their energy signs are opposite because the reduced kinetic form is $\diag(-6,b)$.
Nonlinear boundedness does not follow.
The repair below preserves this entire homogeneous model, including its previously observed nonlinear escape trajectories.

\section{The constrained diagonal scalar representation}
\label{sec:scalar}

Use dimensionless time $\tau$ and divide the action by $a_0^2/\kappa$.
Let $f$ be a real unit-normalized scalar harmonic:
\begin{equation}
 \Delta f=-Lf,\quad L=\ell(\ell+2),\quad
 \int_{S^3}f^2\dd V_\gamma=1,\qquad\omega=\sqrt3.
\end{equation}
The closed scalar representation is
\begin{equation}
 \frac{\delta(X+iZ)}{r_0}=e^{i\omega\tau}
 \{vfY+u\dd Y(\nabla f)+i[zfY+q\dd Y(\nabla f)]\}.
 \label{eq:basis}
\end{equation}
Direct differentiation of the embedding gives
\begin{align}
 \Delta(fY)&=-(L+3)fY+2\dd Y(\nabla f),\\
 \Delta[\dd Y(\nabla f)]&=2LfY-(L-1)\dd Y(\nabla f).
\end{align}
Thus independent scalar-field copies with a common scalar harmonic would not form the correct closed representation.
The integrated Bochner identity with $\operatorname{Ric}_\gamma=2\gamma$ gives
\begin{equation}
 \int|\nabla f|^2=L,\qquad \int|\nabla\nabla f|^2=L(L-2).
 \label{eq:bochner}
\end{equation}
Use clock gauge $\delta\chi=0$ and, for $\ell\ge2$, spatial gauge with zero scalar metric shear:
\begin{equation}
 h_{ij}=a_0^2e^{2\psi f}\gamma_{ij},\qquad N=e^{Af},\qquad N^i=B\nabla^if.
\end{equation}
Put $x=(v,u,z,q,\psi)^T$.
In dimensionless units the complete integrated quadratic Lagrangian is $\mathcal L_g+\mathcal L_k+\mathcal L_s$, where
\begin{align}
 \mathcal L_g={}&-3\psi'^2-2L\psi'B-LB^2+\tfrac32(A+\psi)^2
 +2L\psi(A+\psi)-L\psi^2-\tfrac\Lambda2(A+3\psi)^2,\label{eq:Lg}\\
 \mathcal L_k={}&\tfrac b2\{(v'-\omega z)^2+(z'+\omega v)^2
 +L[(u'-\omega q-B)^2+(q'+\omega u)^2]-2\omega BL(z-q)\}\nonumber\\
 &+b\omega(3\psi-A)(z'+\omega v)+\tfrac{3b+\alpha}{4}(3\psi-A)^2,\label{eq:Lk}\\
 \mathcal L_s={}&-\tfrac b2\{(L+3)(v^2+z^2)-4L(vu+zq)+L(L-1)(u^2+q^2)\}
 \nonumber\\&-b(A+\psi)(3v-Lu)-\tfrac{3b}{4}(A+\psi)^2-\tfrac{bm}{2}v^2.
 \label{eq:Ls}
\end{align}
Here $\Lambda=(6+\alpha)/4$.
The last term of~\eqref{eq:Lk} includes the clock kinetic density.
The shift term $-2\omega BL(z-q)$ is essential to the physical momentum constraint.
Variation of $A,B$ gives
\begin{align}
 A&=-\frac{L(2-\alpha)u+(8L-3\alpha-18)\psi+\sqrt3(\alpha-2)z'+6(\alpha-2)v}{\alpha+6},\label{eq:A}\\
 B&=\frac{(\alpha-2)u'-2\sqrt3(\alpha-2)q+\sqrt3(\alpha-2)z-8\psi'}{\alpha+6}.
 \label{eq:B}
\end{align}
Substitution yields
\begin{equation}
 \mathcal L_{\rm red}=\tfrac12x'^TKx'+x'^TDx-\tfrac12x^THx,
 \qquad Kx''+Gx'+Hx=0,\quad G=D-D^T.
 \label{eq:reduced}
\end{equation}
The explicit kinetic form is
\begin{equation}
 K=\begin{pmatrix}
 b&0&0&0&0\\0&cL&0&0&cL\\0&0&e&0&0\\0&0&0&bL&0\\
 0&cL&0&0&2(8L-3\alpha-18)/(\alpha+6)
 \end{pmatrix},\quad
 c=\frac{2(2-\alpha)}{\alpha+6},\quad e=\frac{\alpha(2-\alpha)}{\alpha+6}.
 \label{eq:K}
\end{equation}
Equations~\eqref{eq:Lg}--\eqref{eq:B} give an explicit algebraic recipe for all entries; Appendix~\ref{app:matrices} also lists $G,H$.

\subsection{Direct field equations and exceptional gauge reduction}

Direct covariant scalar variation gives
\begin{align}
 0&=v''-2\omega z'+(L+m)v-2Lu-6(\psi-A),\label{eq:KGv}\\
 0&=u''-2\omega q'+(L-4)u-2v-B'-A-\psi,\label{eq:KGu}\\
 0&=z''+2\omega v'+Lz-2Lq-\omega(A'-3\psi'-LB),\label{eq:KGz}\\
 0&=q''+2\omega u'+(L-4)q-2z-2\omega B.\label{eq:KGq}
\end{align}
The omitted clock equation is $A'-3\psi'-LB=0$.
Substitution of~\eqref{eq:A}, \eqref{eq:B} and $x''=-K^{-1}(Gx'+Hx)$ gives five exact symbolic zeros.
This checks the coupled physical equations rather than a test-field model.

For $\ell=1$, $L=3$ and $\nabla_i\nabla_jf=-f\gamma_{ij}$, so scalar metric shear is identically absent.
The residual spatial diffeomorphism acts as $u\mapsto u-\xi$, $\psi\mapsto\psi+\xi$.
Its vector $n=(0,-1,0,0,1)^T$ satisfies $Kn=Gn=Hn=0$.
Fix $\psi=0$ and use the four physical coordinates $(v,u,z,q)$ with the principal four-by-four matrices.
Their four accelerations substituted into all five reduced equations give exact zero, including the omitted fifth equation.
No singular kinetic matrix is inverted or physical equation discarded.

The generic physical count is five scalar degrees, four transverse-vector degrees and two tensor degrees: eleven, agreeing with nine scalars and two gravitational polarizations.
Exceptional low harmonics have smaller representation spaces and are treated explicitly.

\section{Instabilities and a radial-potential obstruction}
\label{sec:obstruction}

At~\eqref{eq:radialpoint}, the original $\ell=1$ sector has an unstable quartet
\begin{equation}
 \sigma=\pm1.1175696556\pm1.4428277436i
\end{equation}
and elliptic frequencies $1.7028084808$, $4.0130865119$ in dimensionless time.
At $\ell=3$ it has a real growing/decaying pair $\sigma=\pm0.1975187692$ and four elliptic pairs.
Physical exponents are divided by $a_0$.
Exact Sturm counts of the polynomials in Appendix~\ref{app:roots} establish their root types.
The $\ell=2$ sector has four elliptic pairs and an exact zero vector
\begin{equation}
 x_0=(0,0,1,1/2,0)^T,\qquad Hx_0=0.
\end{equation}
This is a physical internal $U(4)$ orientation variation: an imaginary symmetric traceless target generator produces $\delta\Phi=iS\Phi_0$ with precisely this radial/tangent ratio.
It is not a diffeomorphism.
The stiffness has rank four and the zero characteristic root has multiplicity two, so the unrestricted theory has a secular conjugate direction.
For $\ell\ge4$, exact leading minors expanded at $L=24+t$ prove $K,H>0$ in the original model.

\begin{proposition}
Every parameter pair in~\eqref{eq:windowa}--\eqref{eq:windowm} has a growing physical $\ell=3$ scalar mode in the canonical model.
\end{proposition}
\begin{proof}
The exact stiffness determinant is
\begin{align}
 \det H={}&\frac{L^3\alpha(L-8)(L-3)(\alpha-2)^4(2L-\alpha-6)}{16(\alpha+6)^2}\mathcal F(L),\\
 \mathcal F(L)={}&2L^2-L\alpha+2Lm-14L-\alpha m+24\alpha-6m-48.
\end{align}
At $L=15$ every prefactor is positive, whereas
\begin{equation}
 \mathcal F(15)=192+9\alpha+(24-\alpha)m
 <\frac{3\alpha(11\alpha-94)}{6+\alpha}<0.
\end{equation}
The first inequality uses~\eqref{eq:windowm} and $24-\alpha>0$.
Since $K>0$ for $L>3$, the real continuous determinant $\det(\sigma^2K+\sigma G+H)$ is negative at zero and positive for large positive $\sigma$.
It has a positive zero, proving a growing solution.
\end{proof}

A new radial mass scan cannot solve this obstruction.
Positive target alignment alone changes only the imaginary tangent stiffness and leaves the negative real-block determinant at $\ell=3$ unchanged; it also destabilizes the gyroscopically protected $\ell=2$ sector when used alone.
Traceless spatial strain alone misses $\ell=1$.
These failures motivate the combined action below.

\section{A local preferred-clock repair}

On the timelike-clock patch define
\begin{equation}
 n_\mu=-\frac{\partial_\mu\chi}{\sqrt{-\nabla_\rho\chi\nabla^\rho\chi}},\qquad
 h_{\mu\nu}=g_{\mu\nu}+n_\mu n_\nu.
\end{equation}
Let $R_3$ be the intrinsic curvature of a clock leaf and define real field-space contractions
\begin{align}
 \mathcal B_{\mu\nu}&=h_\mu{}^\rho h_\nu{}^\sigma\sum_I\partial_\rho\Phi^I\partial_\sigma\Phi^I,
 &\mathcal T&=h^{\mu\nu}\mathcal B_{\mu\nu},\\
 \mathcal E_{\mu\nu}&=\mathcal B_{\mu\nu}-\tfrac13\mathcal T h_{\mu\nu},
 &\mathcal Q&=\mathcal T-\tfrac12sR_3,\\
 \mathcal D&=|X|^2|Z|^2-(X\cdot Z)^2.&&
\end{align}
The repaired action is
\begin{equation}
 S=S_0-\int\sqrt{-g}\left[\lambda\mathcal E_{\mu\nu}\mathcal E^{\mu\nu}
 +\zeta\mathcal Q^2+\frac{h_*}{2a_0^2r_0^2}\mathcal D\right]\dd^4x.
 \label{eq:newaction}
\end{equation}
Here $[\lambda]=[\zeta]=L^4$, $h_*$ is dimensionless, and the chosen coefficients are
\begin{equation}
 \bar\lambda=\frac\lambda{\kappa a_0^2}=\tfrac1{10},\qquad
 \bar\zeta=\frac\zeta{\kappa a_0^2}=1,\qquad h_*=6.
 \label{eq:repairpoint}
\end{equation}
These values define the model; no naturalness or ultraviolet matching is assumed.
The Gram determinant $\mathcal D$ is nonnegative, preserving the bounded-below scalar potential.

For every homogeneous ansatz $g=-N^2\dd t^2+a^2\gamma$, $X+iZ=re^{i\theta}Y$, with homogeneous timelike clock,
\begin{equation}
 \mathcal B_{ij}=r^2\gamma_{ij},\quad \mathcal T=3r^2/a^2,\quad R_3=6/a^2,
 \qquad\mathcal E=\mathcal Q=\mathcal D=0.
\end{equation}
All additions and their first variations vanish, preserving the exact background, embedding and full homogeneous action.
This preserves the homogeneous limitations as well.

In the scalar gauge, geometric variation gives
\begin{align}
 \delta\mathcal B_{ij}&=2b[vf\gamma_{ij}+u\nabla_i\nabla_jf],&\delta s&=2bvf,\\
 \delta\mathcal T&=2b(3v-Lu-3\psi)f,&\delta R_3&=4(L-3)\psi f.
\end{align}
Thus
\begin{equation}
 \mathcal E^{(1)}_{ij}=2bu(\nabla_i\nabla_jf-\gamma_{ij}\Delta f/3),\qquad
 \mathcal Q^{(1)}=-2bL(u+\psi)f.
\end{equation}
The added quadratic potential is
\begin{equation}
 W=\frac83\bar\lambda b^2L(L-3)u^2+4\bar\zeta b^2L^2(u+\psi)^2+\tfrac12bLh_*q^2.
 \label{eq:W}
\end{equation}
The factor in the strain term follows from~\eqref{eq:bochner}.
No lapse or shift appears, so $A,B,K,G$ are unchanged and $H_{\rm r}=H+\nabla_x\nabla_xW$.
The new $u,q$ equations acquire $\partial_uW/(bL)$ and $\partial_qW/(bL)$, respectively; the radial equations retain their displayed form.
All four direct repaired scalar residuals and the omitted clock equation again vanish identically.
The residual gauge null vector at $L=3$ is preserved because $u+\psi$ is invariant and $L(L-3)=0$; its dropped equation remains an exact consequence of the four physical equations.

\section{All-harmonic positive energy}
\label{sec:positive}

\begin{theorem}
\label{thm:positive}
At~\eqref{eq:radialpoint}, \eqref{eq:repairpoint}, every physical inhomogeneous scalar and transverse-vector harmonic has strictly positive quadratic kinetic and stiffness forms after constraint elimination and gauge reduction.
The tensor correction is positive.
Every such harmonic has nonzero real normal frequencies and bounded linear motion in its positive energy norm.
This does not claim positive energy or nonlinear stability for the unchanged homogeneous sector.
\end{theorem}

\subsection{Exact scalar proof}

For $L\ge8$, the five leading kinetic minors are
\begin{equation}
 \frac{19}{40},\quad\frac{361L}{1220},\quad\frac{6859L}{744200},\quad
 \frac{130321L^2}{29768000},\quad\frac{130321L^2(L-3)}{14884000}.
\end{equation}
They are positive.
The repaired stiffness splits into real $(v,u,\psi)$ and imaginary $(z,q)$ blocks.
The leading real-block minors are
\begin{align}
 R_1&=\frac{19(1525L+4473)}{61000},\\
 R_2&=\frac{361L(613600L^2-456633L-1596861)}{183000000},\\
 R_3&=\frac{361L(L-3)}{36600000000}
 (259509975L^3-1960056193L^2\nonumber\\
 &\hspace{40mm}+3767182800L-534407580).
\end{align}
For the imaginary block,
\begin{equation}
 I_1=\frac{19L}{610},\qquad I_2=\frac{361L^2(L-2)}{24400}.
\end{equation}
Every displayed minor has positive coefficients and a positive constant term after expansion in $t=L-8$.
This exact rational coefficient certificate proves positivity for every $L\ge8$, not just finitely many harmonics.
The executable saves all expanded coefficients.
At exceptional $L=3$, the four physical stiffness minors are
\begin{equation}
 \frac{21489}{7625},\quad\frac{23064651}{1525000},\quad
 \frac{1314685107}{930250000},\quad\frac{74937051099}{37210000000},
\end{equation}
and the kinetic minors are
\begin{equation}
 \frac{19}{40},\quad\frac{1083}{1220},\quad\frac{20577}{744200},\quad
 \frac{1172889}{29768000}.
\end{equation}
Sylvester's criterion completes the scalar proof.
The conserved positive energy is
\begin{equation}
 E_2=\tfrac12x'^TKx'+\tfrac12x^TH_{\rm r}x,\qquad
 E_2'=-x'^TGx'=0.
 \label{eq:energy}
\end{equation}
Selected frequencies are listed in Table~\ref{tab:freq}.

\begin{table}[htbp]
\centering
\caption{Repaired positive scalar frequencies in units $a_0^{-1}$.}
\label{tab:freq}
\begin{tabular}{cl}\toprule
$\ell$&$a_0\omega_j$\\\midrule
1&0.35222722, 1.66575944, 2.34696391, 5.19186305\\
2&1.15713293, 2.07161365, 2.82603344, 3.60667197, 7.04429409\\
3&2.35771272, 3.25096940, 3.90010792, 4.73791706, 8.91721306\\\bottomrule
\end{tabular}
\end{table}

\subsection{Transverse vectors and exceptional Killing harmonics}

For a unit divergence-free harmonic $T_i$ with rough Laplace eigenvalue $1-n(n+2)$, write $\delta\Phi/r_0=e^{i\omega\tau}\dd Y[(V+iZ)T]$.
Fix the spatial vector metric and retain transverse shift $BT$.
With $d=n(n+2)-3$, the original quadratic Lagrangian is
\begin{equation}
 \mathcal L_V=\frac b2[V'^2+Z'^2+2\omega(VZ'-ZV')-d(V^2+Z^2)]
 -bB(V'-2\omega Z)+\frac12(b+d/2)B^2.
\end{equation}
The gravitational contribution is $dB^2/4$, since $\int|\nabla_{(i}T_{j)}|^2=d/2$.
The shift is $B=2b(V'-2\omega Z)/(d+2b)$.
The repair adds $-2\bar\lambda b^2dV^2-bh_*Z^2/2$; compression has zero transverse first variation.
With $f_d=d/(d+2b)$, the reduced matrices are
\begin{align}
 K_V&=b\diag(f_d,1),&G_V&=\begin{pmatrix}0&-2b\sqrt3f_d\\2b\sqrt3f_d&0\end{pmatrix},\\
 H_V&=b\diag\left(d(1+4\bar\lambda b),\ d+\frac{24b}{d+2b}+h_*\right).
\end{align}
These kinetic and stiffness entries are positive for every $n\ge2$.
At $n=1$, $V$ is a spatial Killing diffeomorphism and disappears from the action; the remaining relative-rotation oscillator has frequency $\sqrt{12+h_*}/a_0=\sqrt{18}/a_0$.
Thus internal target directions are explicitly distinguished from gauge directions.

\subsection{Tensor correction}

For a unit-normalized transverse-traceless metric amplitude, the original closed sector has $\delta\Phi=0$ and
\begin{equation}
 \mathcal L_T=\tfrac18[h_T'^2-(n(n+2)+2b)h_T^2],\qquad n\ge2.
\end{equation}
The tensor rough Laplace eigenvalue is $2-n(n+2)$~\cite{harmonics}.
Although $\delta\mathcal B_{ij}=0$, metric trace removal gives $\mathcal E^{(1)}_{ij}=-bh_{T,ij}$.
The compression and alignment variations vanish.
Therefore $-\bar\lambda b^2h_T^2$ adds $8\bar\lambda b^2=0.1805$ to the squared frequency:
\begin{equation}
 \omega_{T,n}^2=\frac{n(n+2)+2b+8\bar\lambda b^2}{a_0^2}
 =\frac{n(n+2)+2261/2000}{a_0^2}>0.
\end{equation}
Its scalar source vanishes by the tensor trace and divergence identities, so this remains a closed sector.
This completes Theorem~\ref{thm:positive}.

\section{A finite-band stationary state and coupled evolution}

Choose scalar harmonics $\ell=1,2,3$ with all degeneracies $(\ell+1)^2$.
The band contains $4\times4+5\times9+5\times16=141$ canonical oscillators.
It is a finite truncation of the quadratic theory.
Its physical effective-theory interpretation assumes its wavelengths and frequencies are below a suitable cutoff and omitted operators are small; no ultraviolet matching or truncation-error bound is available.

Drop a total derivative to set $D=G/2$, giving $p=Kx'+Dx$ and positive canonical Hamiltonian matrix
\begin{equation}
 \HH=\begin{pmatrix}H_{\rm r}+D^TK^{-1}D&-D^TK^{-1}\\-K^{-1}D&K^{-1}\end{pmatrix}.
\end{equation}
The normalized physical coordinates are
\begin{equation}
 Q=\frac{a_0}{\sqrt\kappa}x,\qquad
 P=\frac{a_0}{\sqrt\kappa}(Kx'+Gx/2),\qquad [Q_i,P_j]=i\delta_{ij}.
 \label{eq:normalization}
\end{equation}
There is no sphere-volume factor because harmonics have unit $L^2$ norm.
For the standard symplectic matrix $J$, a centered pure ground-state covariance is
\begin{equation}
 C=\tfrac12\HH^{-1/2}
 \sqrt{-(\HH^{1/2}J\HH^{1/2})^2}\,\HH^{-1/2}.
 \label{eq:cov}
\end{equation}
Principal positive square roots are used.
The finite-dimensional symplectic normal-form construction~\cite{williamson}, or directly diagonalizing the real antisymmetric matrix $\HH^{1/2}J\HH^{1/2}$, proves
\begin{equation}
 C+iJ/2\ge0,\qquad CJC=J/4,\qquad (J\HH)C+C(J\HH)^T=0.
\end{equation}
The product over harmonic degeneracies is a stationary pure Gaussian.
Its physical coordinate covariance is $(\kappa/a_0^2)C_{QQ}$, so the dimensionless root-mean-square amplitudes scale as $\sqrt\kappa/a_0$.
No nonlinear containment follows from that scaling.

Actual canonical matrices for each multiplet are supplied with the calculation.
Their ordering is all $Q$ coordinates followed by all $P$ coordinates, with $d=4$ for $\ell=1$ and $d=5$ for $\ell=2,3$.
Their stationarity residual norms are $7.15\times10^{-14}$, $2.44\times10^{-13}$ and $4.94\times10^{-13}$.
For $\ell=1$, purity residual is $4.32\times10^{-15}$; the smallest uncertainty eigenvalue is $-2.94\times10^{-15}$, consistent with exact zero uncertainty eigenvalues and numerical roundoff.
Analytic formula~\eqref{eq:cov} establishes positivity, rather than a roundoff threshold.
Finite-band kernels are smooth and lack continuum ultraviolet singularities; they are not called continuum Hadamard states.

For a coupled evolution, take one $\ell=1$ harmonic with $x(0)=(10^{-4},0,0,0)$ and $x'(0)=0$ on the $\psi=0$ gauge slice.
Matrix exponentials propagate the constrained linear system on $0\le\tau\le100$, and~\eqref{eq:A}--\eqref{eq:B} reconstruct the metric.
One obtains
\begin{equation}
 E_2(0)=1.4091147541\times10^{-8},\qquad
 \max\frac{|E_2(\tau)-E_2(0)|}{E_2(0)}=1.32\times10^{-12}.
\end{equation}
Positive energy gives the analytic all-time bound $\|(x,x')\|\le9.51210\times10^{-4}$; the maximum sampled norm is $7.76142\times10^{-4}$.
The maximum reconstructed lapse and shift harmonic amplitudes are $3.50155\times10^{-4}$ and $2.11911\times10^{-4}$.
The metric responds even though spatial scalar metric amplitude is gauge-fixed.
This is a coupled linear metric--matter evolution, not finite-amplitude nonlinear constraint-solving or nonlinear partial-differential-equation evolution.

\section{Principal cones and effective-theory limitations}

The compression term contains $R_3^2$ and fourth spatial derivatives.
In clock gauge, the new operators depend on spatial geometry and spatial matter derivatives; lapse is multiplicative and shift absent from their quadratic potential.
No extra time derivatives occur at this order.
Covariantly, derivatives of the clock enter the foliation geometry; the verified local clock-gauge derivative count is not a global nonlinear Hamiltonian-degeneracy theorem or a result across a null clock gradient.

The constrained high-frequency symbol is more informative than unreduced derivative counting.
The $(u,\psi)$ kinetic block grows as $L$ while its stiffness grows as $L^2$.
The asymptotic matrices are
\begin{equation}
 K_\infty=\begin{pmatrix}c&c\\c&c+2\end{pmatrix},\qquad
 H_\infty=K_\infty+\tfrac{16}{3}\bar\lambda b^2\begin{pmatrix}1&0\\0&0\end{pmatrix}
 +8\bar\zeta b^2\begin{pmatrix}1&1\\1&1\end{pmatrix}.
\end{equation}
Three scalar speeds tend to one.
For the other two, $z=c_s^2$ obeys
\begin{equation}
 2400000z^2-12362000z+10380399=0,
 \qquad z=\frac{6181}{2400}\pm\frac{19\sqrt{920485}}{12000}.
\end{equation}
Thus $c_{s,1}=1.0277822886$ and $c_{s,2}=2.0234863233$.
The reduced dispersion is $\omega^2\sim c_s^2L/a_0^2$, rather than a demonstrated $L^2/a_0^2$ frequency law.
The real vector branch tends to $\sqrt{1+4\bar\lambda b}=\sqrt{1.19}$ and the imaginary branch to one.
The preferred-clock cones differ from the metric null cone; energy positivity does not prove metric microcausality or a Lorentz-invariant ultraviolet completion.

All-harmonic positivity is a mathematical property of the specified quadratic model, not permission to extrapolate an effective theory to arbitrary physical momentum.
An inherited physical Hadamard construction for the original canonical normally hyperbolic system remains a result for that original system.
It is not transferred to~\eqref{eq:newaction}, whose principal symbol changes.
Continuum preferred-clock states and renormalization need separate analysis.
The finite-band state and bounded linear evolution are usable within their stated cutoff and perturbative assumptions.
Interactions with the unchanged homogeneous opposite-energy modes remain the main dynamical limitation: cubic couplings, nonlinear constraints and energy-transfer time bounds are the next substantive calculations.

\section*{Computation, chronology and attribution}

The accompanying executable \texttt{construct\_inhomogeneous.py} supplies the full quadratic forms, exact positivity coefficients, low-mode Sturm certificates, numerical covariance matrices and the reconstructed linear trajectory.
It uses local Python with SymPy 1.14.0 and SciPy 1.17.1.
Symbolic identities establish the sign and constraint claims; numerical work supplies displayed frequencies, covariances and trajectory samples.
The canonical reference and radial homogeneous construction precede this inhomogeneous calculation, performed on 29 October 2026 Beijing time.
Historical work, this local research date and any public-disclosure date are distinct.

This is AI-assisted work without human expert review.
As reported by the operator, earlier work was Fable 5-led; from July 2026 onward results were predominantly produced by GPT models, with DeepSeek and multiple GPT models also participating in cross-review and preparation.
Human authors and affiliations are blank.
No global novelty or priority determination is made.
The attribution uncertainty is retained explicitly:
\begin{quote}
``I genuinely do not know how authorship, intellectual credit, or ownership should be attributed for these results.''
\end{quote}

\appendix
\section{Explicit reduced matrices}
\label{app:matrices}

The only nonzero upper-triangular gyroscopic entries are
\begin{equation}
 G_{13}=\frac{2\sqrt3\alpha(\alpha-2)}{\alpha+6},\quad
 G_{24}=\frac{4\sqrt3L(\alpha-2)}{\alpha+6},\quad
 G_{45}=\frac{4\sqrt3L(2-\alpha)}{\alpha+6}.
\end{equation}
Antisymmetry supplies the lower triangle.
Set $d_0=\alpha+6$.
The original symmetric stiffness has entries
\begin{align}
 H_{11}={}&[-L\alpha^2/4-L\alpha+3L-\alpha^2m/4+9\alpha^2
 -\alpha m-36\alpha+3m+36]/d_0,\\
 H_{12}={}&L(-\alpha^2+8\alpha-12)/d_0,\\
 H_{15}={}&3(4L\alpha-8L-\alpha^2-4\alpha+12)/d_0,\\
 H_{22}={}&L(-2L\alpha+4L+\alpha^2+4\alpha-12)/d_0,\\
 H_{25}={}&H_{22},\\
 H_{55}={}&(16L^2-14L\alpha-84L+3\alpha^2+36\alpha+108)/d_0,\\
 H_{33}={}&L\alpha(2-\alpha)/d_0,\\
 H_{34}={}&2L\alpha(\alpha-2)/d_0,\\
 H_{44}={}&L(-L\alpha^2-4L\alpha+12L-8\alpha^2+64\alpha-96)/(4d_0).
\end{align}
Symmetry supplies opposite entries and all unlisted entries vanish.
The repair adds
\begin{align}
 \Delta H_{22}&=\tfrac{16}{3}\bar\lambda b^2L(L-3)+8\bar\zeta b^2L^2,\\
 \Delta H_{25}&=\Delta H_{52}=\Delta H_{55}=8\bar\zeta b^2L^2,\\
 \Delta H_{44}&=bLh_*.
\end{align}
At $L=3$ delete the fifth row and column after the residual gauge fixing.
At $L\ge8$ retain all five variables.

\section{Exact original low-mode root certificates}
\label{app:roots}

For $w=\omega^2$, the monic characteristic polynomials of the physical original sectors at~\eqref{eq:radialpoint} are
\begin{align}
 p_1(w)={}&w^4-\tfrac{2067}{100}w^3+\tfrac{22361}{250}w^2
 -\tfrac{144303}{500}w+\tfrac{12951}{25},\\
 p_2(w)={}&w^5-\tfrac{2281}{50}w^4+\tfrac{1181621}{2000}w^3
 -\tfrac{26396913}{10000}w^2+\tfrac{3263619}{1250}w,\\
 p_3(w)={}&w^5-\tfrac{4031}{50}w^4+\tfrac{4268341}{2000}w^3
 -\tfrac{112729759}{5000}w^2\nonumber\\
 &+\tfrac{793930209}{10000}w+\tfrac{391482}{125}.
\end{align}
Exact Sturm counts give two positive real roots and a nonreal pair for $p_1$, four positive roots and one zero for $p_2$, and four positive roots and one negative root for $p_3$.
The associated unstable exponents are therefore physical constrained modes.

\bibliographystyle{unsrt}
\bibliography{references}
\end{document}
